\chap{综合案例：一次跨模块改进如何完成}
\label{sec:case}
\subsection{案例条件与初始问题}
本节构造一个以 Portal 类任务平台为背景的设计示例。场景中，用户提交长时间任务，可以查看进度、请求取消并访问结果。平台包含用户界面、任务状态服务、执行器和结果存储。该案例用于具体说明框架中的对象与过程，不表示一次真实生产故障的复盘，也不提供实验效果数字。

初始问题是：用户点击取消后，界面立即显示“已取消”，但执行器仍可能进行清理，资源也尚未释放。用户无法判断是否需要再次操作，更不清楚已有结果是否保留。虽然取消请求已经被成功接收，界面的表达却使用户形成了超出实际证据的理解。

人的初始需求可以表述为“取消后让我知道任务究竟停了没有，已有结果尽量保留”。Leader 将其联系到当前工作，识别出三个需要处理的部分：状态语义、实际停止和结果可见性。前两项是当前问题的核心，输出保留则需要明确何种中间产物能够被解释和继续使用。

\subsection{形成语义约束与工作范围}
团队先用正常、延迟、重复和并发四种情境澄清要求。正常情境下，用户请求取消并最终得到停止确认；延迟情境下，界面显示请求已接收且正在等待；重复情境下，多次请求不会造成相互矛盾的状态；并发情境下，如果任务已经自然完成，界面应显示实际完成结果并解释取消未改变终态。

输出保留的第一版约束可以是：已经正式写入并符合访问条件的结果继续可见，尚未完成的临时数据不作为完整结果展示。这个约束允许先解决状态误导，再由真实反馈决定是否扩展中间结果能力。未确定的部分被明确保留，而不会隐藏在“尽量”这样的模糊词中。

\begin{table}[htbp]\centering\small
\caption{构造案例中的状态含义与可见行为}\label{tab:states}
\begin{tabularx}{\linewidth}{p{2.5cm}YY}\toprule
状态或事实 & 由何种事实支持 & 用户可以理解和采取的行动\\\midrule
运行中 & 执行器仍执行任务 & 查看进度或提交取消请求\\
取消请求已接收 & 请求已持久记录并进入处理 & 知道操作生效，继续等待确认\\
停止已确认 & 执行器报告停止，必要清理条件成立 & 查看终态和可保留结果\\
自然完成 & 完成事件已确定实际终态 & 查看结果，并理解取消发生得更晚\\
停止处理异常 & 实际停止尚未得到可靠确认 & 查看异常说明和可用处理入口\\\bottomrule
\end{tabularx}\end{table}

任务负责人据此组织一个跨模块改进。状态服务成员负责请求幂等与终态语义，执行器成员负责停止确认与资源观察，界面成员负责用户可见表达，验证成员核对完整任务。长期职责保持清楚，临时负责人则跟踪跨边界协议与实际交付。

\subsection{协议设计与并发条件}
协议区分“请求已接收”与“停止已确认”。前者表示系统已经记录用户意图，后者表示执行结果满足终态条件。执行器延迟不应被界面误认为操作丢失，因此状态服务需要提供可持续查询或订阅的中间状态。重复请求应返回一致的当前事实。

取消与自然完成可能竞争。若完成事件先确立了终态，之后的取消请求不应覆盖真实结果。若取消先进入处理，任务仍可能在停止信号生效前完成。具体规则可以由状态版本、事件顺序和执行器确认共同支持，但共同目标是让终态解释与实际行为一致。

\fig{13-case-sequence}{构造案例中的取消协议与并发分支。请求接收、执行器停止和界面终态分别由不同事实支持；完成事件竞争时，需要核对实际终态。}{case-sequence}

协议解释通过时序图进入 Workbench。人可以调整停止延迟，观察为什么出现等待状态；也可以让完成事件提前到达，查看终态说明。这种演示帮助人判断行为是否符合期望，并让模块成员以相同例子核对各自实现。

\subsection{并行实现、检查与局部交付}
状态服务与执行器可以并行修改，在协议边界上核对共同情境。界面则可以先使用模拟事件完成状态呈现，再连接实际服务。需要交付的批次包含状态语义和相应消费方；与此无关的页面整理或执行优化不必进入同一批次。

代码检查覆盖重复请求、终态不被错误覆盖以及停止确认的条件。TAU 从普通用户入口创建任务、提出取消、等待确认并访问结果。后台资源检查作为补充证据，核对实际停止与释放条件。两类证据分别说明内部规则与用户体验，彼此可以相互解释。

发布工具捕获实际内容并合成到当前运行版本，记录本次变化与验证范围。其他成员在此期间继续开发无关功能。若后续修改只影响任务列表的布局，已有停止逻辑证据仍可保留；若修改状态映射，则需要更新相关用户旅程检查。

\subsection{一次失败怎样改变设计}
在构造的首轮体验中，假设 Agent 能完成取消，但发现页面显示“取消中”之后没有说明等待什么，也没有提供结果入口。它通过刷新页面才看见终态。这意味着后端协议已经改善，用户仍然缺乏过程反馈。验证报告应记录这种摩擦，并关联到具体界面状态。

任务负责人将问题交回界面职责，补充等待说明和状态更新方式。若结果保留已经由服务支持，则把结果入口放在终态旁边。之后重跑相关体验和状态映射检查，无需重复检查所有无关业务功能。修复通过同一局部交付路径送达。

\fig{14-case-iteration}{构造案例中的迭代过程。验证发现表达问题后，修复回到相关职责；人体验新结果并提出下一项需求，仍有效的工作和证据继续复用。}{case-iteration}

人体验新版 Workbench 后，可能提出“取消后希望看到已经产生的日志片段”。这是一项新的结果保留需求，而不只是界面提示调整。团队据此扩展结果契约，重新判断访问、完整性和存储条件。原有取消请求与终态验证仍可复用，新增检查集中在日志保留及其呈现。

\subsection{案例留下的可复用资产}
本轮工作可以形成几类资产：明确的取消协议、状态转移检查、可重复的用户任务情境、解释并发条件的演示，以及关于用户误解的经验。它们分别服务于实现、验证、沟通和后续学习。将这些材料关联到同一行为对象，有助于下一次修改快速建立背景。

同时，需要避免保存过度具体的临时知识。一次运行的全部日志可以保留为证据入口，长期经验则概括何种条件下会误导用户、如何区分状态以及怎样检查。若未来执行器架构改变，经验的用户目标可能仍成立，具体实现步骤则需要重新核对。

案例展示了框架的组合意义：语义约束帮助明确成功含义，架构职责支持并行深入，局部交付缩短反馈时间，TAU 暴露体验问题，Workbench 让人带着具体情境继续引导。每项机制仍有成本，组合后的净收益需要通过下一节的评价方法检验。
